% 02_setting.tex
\section{Setting and the two exact tools}
\label{sec:setting}

\paragraph*{Setting (paper~E).}
We use the setting and notation of paper~E~\cite{paperE} and recall only what is needed. Near the hyperbolic top the pendulum's phase
space is described by coordinates $(v,w)$ along the unstable and stable directions; $z=(v+w)/\sqrt2$ is the angle measured from the top.
A state is a probability density $\rho(v,w)$; the observer sees $p=\rho*\mathcal N(0,\sigma^2I)$. With the score $s=\nabla\ln p$, the
Hessian $H=-\nabla^2\ln p$, the posterior mean $m=y+\sigma^2s(y)$ (Tweedie~\cite{Efron2011}), $\phi=1-\cos m_z$, and the Fisher informations
$\Jvv=\E[s_v^2]$, $\Jww=\E[s_w^2]$, put $a=\sigma^2\Jvv$ and $\mathrm{den}=\sigma^2(\Jvv+\Jww)$. Paper~E's cancellation-free formula gives the
resolution gap (gap per unit blur) as
\begin{equation}
g=\frac{2\sigma^2\Jww+N}{\sigma\,\mathrm{den}},
\label{eq:gap}
\end{equation}
where $N$ is paper~E's curvature term ($N=-\sigma^{-2}\E[\mathrm{Cov}(x_p,x_z-\sin x_z\,|\,y)]$ with $x_p$ the momentum coordinate), and its exact numerator decomposition $\sigma\,\mathrm{den}\,g=2A_w+N_\phi+(B+R_3)$ with $A_w=\sigma^2\E[s_w^2(1-\phi/4)]$ and
$N_\phi=\tfrac{\sigma^2}{2}\E[s_v^2\phi]$. Paper~E's Theorem~2 bounds $g\ge(a/\mathrm{den})\sqrt{P_w}(1-\varepsilon)$ with
$P_w=4\sigma^2\Jww(1-\Phi/4)(1-\delta_w)\E[s_v^2\phi]/a^2$ and the remainder $\varepsilon=-(B+R_3)/(\sigma a\sqrt{P_w})$, a state quantity;
$\Phi=\E[\phi]$, $\delta_w$ and $\delta_c$ are the decorrelation defects defined there. The \emph{strip} family of paper~E is
\begin{equation}
\begin{aligned}\rho={}&\mathcal N(v;L_v^2)\,\mathcal N(w;L_w^2)\\&\times\big(1+c\,e^{-w^2/2\ell^2}\cos(k_0v/\sigma)\big),\quad 0<c\le1,\end{aligned}
\label{eq:strip}
\end{equation}
with $\ell=\lambda\sqrt\sigma$, $L_w=M\ell$, $L_v=\eta\ell$. We write $\ls=\ln(1/\sigma)$.

\paragraph*{Tool 1: the branch identity \proofst.}
Write $(a)=2A_w/(\sigma\,\mathrm{den})$ (the $w$-branch) and $(a')=N_\phi/(\sigma\,\mathrm{den})$ (the $N$-branch). The last term of the numerator
decomposition equals $-\varepsilon\,(a/\mathrm{den})\sqrt{P_w}$, and $(a/\mathrm{den})\sqrt{P_w}=2\sqrt{(a)(a')}\le(a)+(a')$ (paper~E's Theorem~2 is the
AM--GM inequality of the two branches). Hence, for every state,
\begin{equation}
(1-|\varepsilon|)\big((a)+(a')\big)\;\le\;g\;\le\;(1+|\varepsilon|)\big((a)+(a')\big).
\label{eq:branch}
\end{equation}
Both sides of every result below inherit $\varepsilon$ only through this identity.

\paragraph*{Tool 2: the window inequality.}
Let $\keff^2=\sigma^2\int(\partial_v\partial_w\sqrt p)^2/\int(\partial_w\sqrt p)^2$ (the $v$-fineness of the $w$-gradient) and
$C_w=\E[s_v^2\cos(m_z/2)(1-\sigma^2(H_{ww}+H_{vw}))]/\Jvv$.

\begin{theorem}[window inequality] \proofst
\label{thm:AT2}
For every state with sufficient decay,
\begin{equation}
\begin{aligned}g\;&\ge\;\frac{F_c}{2\keff},\\ F_c&:=\frac{a}{\mathrm{den}}(1-\varepsilon)\,|C_w|\sqrt{(1-\Phi/4)(1-\delta_w)} .\end{aligned}
\end{equation}
Moreover $\keff^2=\Lambda/4$ identically, with paper~E's $\Lambda=\sigma^2\E[(s_vs_w-2H_{vw})^2]/\Jww$.
\end{theorem}
\begin{proof}
Put $f=s_v\sqrt p=2\partial_v\sqrt p$ and $\psi=2\sin(m_z/2)$, so $\psi^2=2\phi$, and from $m=y+\sigma^2s$,
$\partial_w\psi=\cos(m_z/2)(1-\sigma^2(H_{ww}+H_{vw}))/\sqrt2$. Integrating by parts in $w$, $\int f^2\partial_w\psi=-2\int\psi f\partial_wf$, and
Cauchy--Schwarz (a Heisenberg-type step~\cite{Stam1959}) gives $\int(\partial_wf)^2\ge C_w^2\Jvv^2/(16\E[s_v^2\phi])$. The mixed derivatives commute,
$\partial_wf=2\partial_v\partial_w\sqrt p$, and $\int(\partial_w\sqrt p)^2=\Jww/4$, so $\int(\partial_wf)^2=\keff^2\Jww/\sigma^2$. Inserting
$\Jww\ge\sigma^2C_w^2\Jvv^2/(16\keff^2\E[s_v^2\phi])$ into $P_w$, the factor $\E[s_v^2\phi]$ cancels:
$P_w\ge C_w^2(1-\Phi/4)(1-\delta_w)/(4\keff^2)$, and paper~E's Theorem~2 gives the claim. The identity $\keff^2=\Lambda/4$ follows from
paper~E's Lemma~3, $s_w\sqrt p=2\partial_w\sqrt p$ and $(s_vs_w-2H_{vw})\sqrt p=4\partial_v\partial_w\sqrt p$.
\end{proof}

The window inequality is the sharp form of paper~E's Corollary~4: the constant $|C_w|\sqrt{(1-\Phi/4)(1-\delta_w)}$ tends to $1$ on the
states considered here, instead of $(\sqrt2/\pi)(1-\varepsilon_m)\sqrt{1-\varepsilon_\pi^2}\approx0.45$. It contains no window width and no
decorrelation defect $\delta_c$: a window that hides $v$-information where the pendulum is straight pays for it with $w$-information in
exactly the inverse amount. What remains is the fineness $\keff$, and the question is how large it can be. On the strips
$C_w=1.00000$ and the inequality is saturated as $\sigma\to0$.
